\section{블루투스 지원 --- Web API와 파일 기반 웹앱 제작 시 요구사항}

본 장은 \code{Web/bluetooth.js}의 \code{BluetoothManager}를 중심으로 Web
Bluetooth 통신 구조를 분석하고, 특히 \code{file://} 환경에서 단독 실행되는
웹앱을 만들 때의 요구사항을 정리한다.

\subsection{연결 흐름}

\code{BluetoothManager.connect()}는 다음 순서로 BLE 연결을 수립한다.

\begin{enumerate}
  \item \textbf{장치 탐색}: \code{navigator.bluetooth.requestDevice()} 호출.
        이름 필터(\code{PicoBLE}, \code{HMSoft}, \code{BT05})와 UART 서비스 필터를
        병행한다.
  \item \textbf{GATT 연결}: \code{device.gatt.connect()} 후 안정화를 위한 지연.
  \item \textbf{서비스/특성 획득}: UART 서비스(\code{0xFFE0})와 특성(\code{0xFFE1})을 획득.
  \item \textbf{알림 구독}: \code{startNotifications()}로 수신 알림을 구독하고,
        실패 시 폴링 모드로 폴백.
  \item \textbf{연결 해제 핸들러 등록}: \code{gattserverdisconnected} 이벤트 처리.
\end{enumerate}

이 과정을 시퀀스로 보면 그림~\ref{fig:bleseq}와 같다. 모든 단계는 사용자의
``연결'' 버튼 클릭(제스처) 안에서 시작되어야 하며, 알림 구독에 실패하면
주기적 읽기(폴링)로 자동 전환된다.

\begin{diagram}{BLE 연결 수립 시퀀스}{fig:bleseq}
\node[dom, text width=34mm] (web) at (0,0) {웹 앱\\{\scriptsize BluetoothManager}};
\node[dom, text width=30mm] (dev) at (8,0) {BLE 장치\\{\scriptsize HM-10 / BT05}};
\draw[rulegray, dashed] (web.south) -- ++(0,-6.2);
\draw[rulegray, dashed] (dev.south) -- ++(0,-6.2);
\coordinate (w) at (web.south);
\coordinate (d) at (dev.south);
\foreach \y/\txt in {-0.7/{requestDevice() — 이름·서비스 필터},
  -1.7/{gatt.connect()},
  -2.7/{getPrimaryService(0xFFE0)},
  -3.7/{getCharacteristic(0xFFE1)},
  -4.7/{startNotifications()}}
  \draw[flow] ([yshift=\y cm]w) -- node[lbl, above]{\scriptsize\txt} ([yshift=\y cm]d);
\draw[back] ([yshift=-5.7cm]d) -- node[lbl, below]{\scriptsize characteristicvaluechanged (수신 알림)} ([yshift=-5.7cm]w);
\end{diagram}

\subsection{청크 송신}

BLE 특성 쓰기는 한 번에 20바이트(\code{MAX\_CHUNK\_SIZE})로 제한되므로,
\code{sendData()}는 명령 문자열을 인코딩한 뒤 20바이트 단위로 분할 전송한다.

\captionof{listing}{청크 송신 핵심 로직}
\begin{minted}{javascript}
for (let i = 0; i < encoded.length; i += MAX_CHUNK_SIZE) {
  const chunk = encoded.slice(i, Math.min(i + MAX_CHUNK_SIZE, encoded.length));
  if (useWithoutResponse)
    await characteristic.writeValueWithoutResponse(chunk); // 빠름(GATT ACK 없음)
  else
    await characteristic.writeValueWithResponse(chunk);
  // 마지막 청크 뒤에는 지연을 두지 않음 -> 단일 청크 명령 지연 제거
  if (i + MAX_CHUNK_SIZE < encoded.length)
    await this.delay(CHUNK_DELAY); // 100ms
}
\end{minted}

청크 간 100\,ms 지연은 HM-10/BT05의 연결 인터벌(약 30\textasciitilde 70\,ms)을
상회하여, 다중 청크 명령(\code{BATCH}, LED 패턴, \code{SYS\_SET})의 패킷 손실을
방지한다. 단일 청크 명령(\code{LED\_ON} 등)은 마지막 청크 뒤 지연을 생략하여
응답성을 높인다. 그림~\ref{fig:chunk}은 긴 명령이 20바이트 청크로 분할\textbf{·}
재조립되는 과정을 보여준다.

\begin{diagram}{20바이트 청크 분할과 펌웨어 측 개행 재조립}{fig:chunk}
\node[boxw, text width=42mm] (cmd) {명령 문자열 + \textbackslash n\\{\scriptsize 예: "BATCH;LED\_ON,0,1|..."}};
\node[hw, below left=8mm and -8mm of cmd, text width=16mm] (c0) {청크0\\{\scriptsize 20B}};
\node[hw, right=4mm of c0, text width=16mm] (c1) {청크1\\{\scriptsize 20B}};
\node[hw, right=4mm of c1, text width=16mm] (c2) {청크2\\{\scriptsize $\le$20B}};
\node[boxw, below=18mm of cmd, text width=42mm] (buf) {피코 수신 버퍼\\{\scriptsize 개행까지 누적(512B)}};
\node[boxw, right=10mm of buf, text width=28mm] (line) {완성 라인\\$\to$ process()};
\draw[flow] (cmd) -- (c0.north);
\draw[flow] (c0) -- node[lbl,right]{100ms} (buf.west);
\draw[flow] (c1) -- (buf.north);
\draw[flow] (c2) -- (buf.east);
\draw[flow] (buf) -- node[lbl,above]{\textbackslash n 검출} (line);
\end{diagram}

\subsection{수신 및 응답 조립}

알림 핸들러는 수신 청크를 누적 버퍼에 모으고 개행 단위로 완성된 메시지를
추출한다. 메시지 종류에 따라 분기 처리한다.

\begin{itemize}
  \item \code{STATUS,...} --- 로버 상태. iframe(대시보드)으로 \code{postMessage} 전달.
  \item \code{SYS\_VALUES,...} --- 시스템 설정 값.
  \item \code{CALIB\_VALUES,...} --- 보정 값.
  \item \code{DIST:}, \code{MAG:} --- 센서 값. 정규식으로 추출하여 Blockly 변수에 반영.
\end{itemize}

응답 대기는 프라미스 기반이다. 빠른 피코 응답과의 경쟁 상태를 피하기 위해,
명령 전송 \emph{이전}에 \code{state.pendingResolve}를 설정하고, 5초
(\code{RESPONSE\_TIMEOUT}) 타임아웃과 연결 해제 시 정리(reject)를 둔다.

\subsection{Web Bluetooth API 요구사항}

\begin{infobox}[title=보안 컨텍스트 및 사용자 제스처]
\begin{itemize}
  \item Web Bluetooth는 \textbf{보안 컨텍스트}에서만 동작한다: HTTPS 또는
        \code{http://localhost}.
  \item \code{requestDevice()}는 \textbf{사용자 제스처}(클릭/터치) 안에서 호출해야
        한다. ARES는 ``신호 연결'' 버튼 클릭으로 이를 만족한다.
  \item Chromium 계열(Chrome, Edge) 데스크톱에서는 \code{file://}도 보안
        컨텍스트로 인정되어 Web Bluetooth가 동작한다. 모바일 Chrome은 더
        엄격할 수 있어 HTTPS 호스팅을 권장한다.
\end{itemize}
\end{infobox}

\subsection{파일 기반(\code{file://}) 웹앱 제작 시 요구사항}

웹 소스를 그대로 \code{file://}로 열면 ES 모듈 임포트가 동작하지 않는다.
\code{<script type="module">}은 \code{file://}에서 \code{origin: null}로 취급되어
동일 출처 정책\footnote{\textbf{동일 출처 정책(Same-Origin Policy)}: 한 출처(origin)의
문서\textbf{·}스크립트가 다른 출처의 리소스에 접근하는 것을 제한하는 브라우저 보안 규칙.
\code{file://}는 출처가 \code{null}이 되어 모듈 \code{import}가 막힌다.}에 의해 \code{import}가 차단되기 때문이다. 따라서 오프라인 단독
실행본을 만들려면 다음 요구사항을 충족해야 한다(상세 빌드 절차는 8장 참조).

\begin{enumerate}
  \item \textbf{ES 모듈 번들링}: \code{esbuild}\footnote{\textbf{esbuild}: 자바스크립트\textbf{·}CSS를
        매우 빠르게 하나로 묶어 주는 번들러(빌드 도구).}로 \code{main.js}와 의존 모듈을
        단일 IIFE\footnote{\textbf{IIFE(Immediately Invoked Function Expression, 즉시 실행 함수
        표현식)}: 정의와 동시에 한 번 실행되는 함수. 내부 변수를 외부와 격리(캡슐화)하는 데 쓴다.}
        번들(\code{main.bundle.js})로 묶어 \code{import} 차단을 제거한다.
  \item \textbf{외부 라이브러리 로컬화}: Blockly, three.js 등을 CDN이 아닌
        \code{vendor/} 로컬 사본으로 동봉한다.
  \item \textbf{스크립트 로딩 순서 보장}: \code{type="module"}을 제거하고
        \code{defer} 스크립트 체인으로 전환하여 DOM 준비 후 실행 순서를 맞춘다.
  \item \textbf{자산 인라인화와 \code{fetch} 심(shim)\footnote{\textbf{심(shim)}: 기존 함수\textbf{·}API
        호출을 가로채 다른 동작으로 대체하거나 보완하는 보정 코드. 여기서는 \code{fetch}를
        가로채 인라인된 자산을 돌려준다.}}: \code{lesson.json},
        예제 XML, GLB 모델 등 \code{fetch}로 로드하던 자산을 base64/텍스트로
        인라인하고, \code{window.fetch}를 가로채는 심을 설치해 네트워크 없이
        자산을 공급한다.
\end{enumerate}

이러한 처리를 통해, 인터넷이 없는 교실 환경에서도 USB나 내부망으로 배포된
단일 폴더를 브라우저로 더블클릭하여 블록 코딩과 BLE 로버 제어를 모두 수행할
수 있다.

\begin{teachbox}{무선 통신의 제약과 우회}
\teachhead{설계 의도}
\begin{itemize}
  \item BLE 20바이트\textbf{·}연결 인터벌 한계에 맞춰 청크 사이에 지연을 둔다.
  \item 응답을 받기 전에 프라미스를 \emph{먼저} 등록해, 빠른 응답과의 경쟁 상태를 피한다.
  \item \code{file://} 제약은 빌드(번들 + 심)로 우회한다.
\end{itemize}
\teachhead{분석할 때 핵심 질문}
\begin{itemize}
  \item 청크 지연을 0으로 하면 긴 명령에서 무엇이 깨질까?
  \item 왜 단일 청크 명령은 마지막 지연을 생략하는가?
\end{itemize}
\teachhead{점검 요소}
\begin{itemize}
  \item \code{CHUNK\_DELAY} 등 타이밍 상수를 바꿔 가며 통신 안정성과 반응 속도의 변화를 관찰하라.
\end{itemize}
\end{teachbox}

